% area: Algebraic Topology

\begin{definition}[Homotopy]
  \label{def:homotopy}
  A homotopy $h : p \simeq q$ between maps $p,q : X \to Y$ is a continuous map 
  \[
    h: X \times I \to Y 
  \]
  such that 
  \[
    h(x,0) = p(x),  h(x,1) = q(x)
  \]
  where $I = [0,1]$ the unit interval.
\end{definition}

\begin{definition}[Equivalence Class/Relation]
  \label{def:eqclass}
An equivalence class is a subset of a larger set containing elements that are considered 
"equivalent" to eachother in the context of a equivalence relation. 

An equivalence relation is a binary operation denoted $a \sim b$ that satisfies three fundamental properties: 

- (Reflexivity) $a \sim a$ 

- (Symmetry) If $a \sim b$, then $b \sim a$ 

- (Transitivity) If $a \sim b$ and $b \sim c$, then $a \sim c$ 

\end{definition}

\begin{example}[An Example of a Homotopy]
  \label{ex:homotopy}
  \ref{def:homotopy}
Let
$$
\gamma_0(t) = (t,0), \quad \gamma_1(t) = (t,t), \quad t \in [0,1].
$$

Both paths go from $(0,0)$ to $(1,1)$.

Define a homotopy $H : [0,1] \times [0,1] \to \mathbb{R}^2$ by
$$
H(t,s) = (t, s t).
$$


Proof that this is infact a homotopy: 


\textbf{Initial path:}
$$
H(t,0) = (t,0) = \gamma_0(t).
$$

\textbf{Final path:}
$$
H(t,1) = (t,t) = \gamma_1(t).
$$

\textbf{Endpoints:}
$$
H(0,s) = (0,0), \quad H(1,s) = (1,s).
$$

Thus, $H$ defines a homotopy between $\gamma_0$ and $\gamma_1$.
\end{example}

\begin{definition}[Loops]
  \label{def:loops}
  A loop is a path 

  \[
    f : I \to X 
  \]
  such that $f(0)=f(1)$
  
\end{definition}

\begin{lemma}[Homotopy Equivalence Class]
  \label{lemma:homeqclass}
  Paths being homotopic defines an equivalence class. Depends on \ref{def:eqclass} and \ref{def:homotopy}
\end{lemma}

\begin{proof}
  \textbf{Reflexivity}
  
Let $\gamma : [0,1] \to X$ be a path. Define
$$
H(t,s) = \gamma(t).
$$
Then $H$ is continuous and satisfies
$$
H(t,0) = \gamma(t), \quad H(t,1) = \gamma(t),
$$
and
$$
H(0,s) = \gamma(0), \quad H(1,s) = \gamma(1).
$$
Thus $\gamma \sim \gamma$.

\textbf{Symmetry}

Suppose $\gamma_0 \sim \gamma_1$ via a homotopy $H(t,s)$. Define
$$
G(t,s) = H(t, 1 - s).
$$
Then $G$ is continuous and satisfies
$$
G(t,0) = H(t,1) = \gamma_1(t), \quad G(t,1) = H(t,0) = \gamma_0(t),
$$
and
$$
G(0,s) = \gamma_0(0), \quad G(1,s) = \gamma_0(1).
$$
Thus $\gamma_1 \sim \gamma_0$.

\textbf{Transitivity}

Suppose $\gamma_0 \sim \gamma_1$ via $H$ and $\gamma_1 \sim \gamma_2$ via $K$. Define
$$
G(t,s) =
\begin{cases}
H(t, 2s) & \text{if } 0 \le s \le \tfrac{1}{2}, \\
K(t, 2s - 1) & \text{if } \tfrac{1}{2} \le s \le 1.
\end{cases}
$$
Then $G$ is continuous (by the pasting lemma), since at $s = \tfrac{1}{2}$ we have
$$
H(t,1) = \gamma_1(t) = K(t,0).
$$
Moreover,
$$
G(t,0) = H(t,0) = \gamma_0(t), \quad G(t,1) = K(t,1) = \gamma_2(t),
$$
and
$$
G(0,s) = \gamma_0(0), \quad G(1,s) = \gamma_0(1).
$$
Thus $\gamma_0 \sim \gamma_2$.
\end{proof}


\begin{definition}[Fundamental Group]
  \label{def:fundgrp}
  \ref{def:loops}
  \ref{def:homotopy}

  \ref{lemma:homeqclass}
  Let $f$ be a loop. 
  Let $[f]$ denote the equivalence class under homotopy for $f$.
  
  We define $\pi_1(X,x)$ to be the set of equivalence classes of loops that start and end at $x$. 

  We will show that this is a group. This is the \textbf{fundamental group} of $X$.
\end{definition}

\begin{lemma}[Fundamental Group Axiom Verification]

  \ref{def:fundgrp}
  The Fundamental Group is indeed a group.

\end{lemma}


\begin{proof}

Let $X$ be a topological space and $x \in X$. The fundamental group $\pi_1(X,x)$ is defined as the set of homotopy classes (relative to endpoints) of loops based at $x$, i.e.
$$
\pi_1(X,x) = \{ [\gamma] \mid \gamma : [0,1] \to X, \ \gamma(0)=\gamma(1)=x \}.
$$

We define a binary operation on $\pi_1(X,x)$ by concatenation:
for loops $\gamma_1, \gamma_2$, define
$$
(\gamma_1 * \gamma_2)(t) =
\begin{cases}
\gamma_1(2t) & 0 \le t \le \tfrac{1}{2}, \\
\gamma_2(2t - 1) & \tfrac{1}{2} \le t \le 1.
\end{cases}
$$
Then define
$$
[\gamma_1] \cdot [\gamma_2] := [\gamma_1 * \gamma_2].
$$

We show this makes $\pi_1(X,x)$ a group.




Suppose $\gamma_1 \sim \gamma_1'$ and $\gamma_2 \sim \gamma_2'$ (homotopy relative endpoints). Then there exist homotopies $H$ and $K$ between them. One constructs a homotopy between $\gamma_1 * \gamma_2$ and $\gamma_1' * \gamma_2'$ by concatenating $H$ and $K$ in the same way:
$$
F(t,s) =
\begin{cases}
H(2t,s) & 0 \le t \le \tfrac{1}{2}, \\
K(2t-1,s) & \tfrac{1}{2} \le t \le 1.
\end{cases}
$$
Thus $[\gamma_1 * \gamma_2] = [\gamma_1' * \gamma_2']$, so the operation is well-defined.

Let $\gamma_1, \gamma_2, \gamma_3$ be loops at $x$. Then
$$
(\gamma_1 * \gamma_2) * \gamma_3 \sim \gamma_1 * (\gamma_2 * \gamma_3),
$$
since both paths traverse $\gamma_1$, $\gamma_2$, and $\gamma_3$ in order, differing only by a reparametrization of $[0,1]$. This reparametrization yields a homotopy relative endpoints. Hence
$$
([\gamma_1] \cdot [\gamma_2]) \cdot [\gamma_3]
= [\gamma_1] \cdot ([\gamma_2] \cdot [\gamma_3]).
$$


Let $e : [0,1] \to X$ be the constant loop at $x$, defined by $e(t)=x$. Then for any loop $\gamma$,
$$
e * \gamma \sim \gamma, \quad \gamma * e \sim \gamma,
$$
by a simple reparametrization homotopy. Hence $[e]$ is the identity element.


Let $\gamma$ be a loop at $x$. Define the reverse path
$$
\gamma^{-1}(t) = \gamma(1 - t).
$$
Then
$$
\gamma * \gamma^{-1} \sim e, \quad \gamma^{-1} * \gamma \sim e,
$$
since the path goes out along $\gamma$ and returns along the same path, which can be contracted to the constant loop. Thus $[\gamma^{-1}]$ is the inverse of $[\gamma]$.

Therefore, the fundamental group is indeed a group.
\end{proof}

\begin{remark}[Dependence on the Base Point]
  \label{rmk:deponbasepoint}
  \ref{def:fundgrp}
  For a path $a : x \to y$, define 
  \[
    \gamma[a] : \pi_1(X,x) \to \pi_1(X,y)
  \]

  by 
  \[
    \gamma[a][f] = [a \cdot f \cdot a^{-1}]
  \]

  This is a homomorphism of groups.
\end{remark}


\begin{problem}
  Prove that the map $\gamma[a]$ as defined in \ref{rmk:deponbasepoint} is indeed a group homomorphism.
\end{problem}


\begin{proposition}
  \label{proposition:continuous-d-s1}
  \ref{def:fundgrp}
  There is no continuous $r: D^2 \to S^1$ such that $r \circ i = \operatorname{id}$. 
\end{proposition}

\begin{proof}
Suppose for contradiction that there does exist such a map $r$. Then the homomorphism 
\[
  \pi_1 (S^1, 1)\xrightarrow{i_*} \pi_1(D^2, 1) \xrightarrow{r_*} \pi_1(S^1, 1)
\]

would be the identity. Since the identity homomorphism of $\mathbb{Z}$ does not through the zero group, this is impossible. 
\end{proof}
